
\subsection{Shen-Ross uniqueness channel across twelve metropolitan markets}
\label{sec:shen-12city}

The Shen-Ross uniqueness channel \citep{shen2021information} measures
how much listing language departs from the corpus average for its
neighborhood, using a Doc2Vec PV-DM embedding trained on the full
listing corpus and quantifying departure as the cosine distance
between each listing and the centroid of the $k$ nearest peer
listings. We follow Shen's recipe verbatim, reproducing her
preprocessing and hyperparameters, and apply the resulting per-listing
uniqueness score as the treatment variable in a partially-linear DML
estimator with cross-fitted ridge nuisances and a Cholesky-solver
ridge solver throughout. Results are reported per standard
deviation of the uniqueness score, with $95\%$ confidence intervals
from a residual bootstrap with $B = 1000$ resamples.

Table~\ref{tab:shen-12city} reports the per-$\sigma$ uniqueness effect
for each of the twelve metropolitan markets in our sample, together
with the bootstrap confidence interval and the binary indicator for
zero exclusion. Two markets identify a Shen-Ross uniqueness effect at
the $95\%$ level. San Francisco recovers a per-$\sigma$ effect of
$+0.131$ with cluster-bootstrap confidence interval
$[+0.055, +0.215]$, within $0.7$ standard errors of the
Atlanta point estimate Shen reports in her 2021 Journal of Urban
Economics study, despite the difference in city. Dallas recovers a
smaller per-$\sigma$ effect of $+0.048$ with confidence interval
$[+0.020, +0.098]$. The remaining ten markets return point estimates
indistinguishable from zero at the $95\%$ level, with confidence
intervals that span both signs and central tendencies clustered near
zero. Atlanta itself, the market on which Shen's identification was
originally established, returns a per-$\sigma$ effect of $-0.006$
with confidence interval $[-0.061, +0.147]$, consistent with the wide
confidence band one would expect at our sample size relative to her
$N \approx 25{,}000$.

Three readings of this pattern are worth distinguishing. The first
attributes the cross-city heterogeneity to power.  Our deduplicated
corpus contains 285 to 334 listings per market, and a Gelman-Carlin
\citep{gelman2014beyond} retrodesign of the per-market estimator
against the assumed true effect of Shen's Atlanta point estimate
($+0.050$ per $\sigma$) finds a median per-market power of $25\%$,
well below the conventional $80\%$ target. Under this reading, the
absence of significance in ten markets is uninformative about the
existence of a Shen-Ross effect; it is exactly what one expects when
power is $25\%$. The two markets where we do detect significance
(San Francisco at $+0.131$, Dallas at $+0.048$) carry Type-S sign-
error rates below $1\%$ each, so the sign of these findings is
reliable, but the magnitude is over-estimated by a Type-M factor of
roughly $2$ on each: the SF point estimate of $+0.131$ is consistent
with a true effect closer to $+0.060$, and Dallas's $+0.048$
consistent with a true effect closer to $+0.025$. We interpret the
two market-level significant findings as honest detections of the
sign, not as unbiased magnitude estimates.

The second reading attributes the cross-city heterogeneity to
genuine between-market variation in the operative residual axis.
Section~\ref{sec:meta-regression-12} addresses this directly with a
random-effects meta-regression that combines the twelve per-market
estimates with Paule-Mandel $\tau^2$ and Hartung-Knapp-Sidik-Jonkman
small-sample CIs.  The primary moderator hypothesis, motivated by
the housing-search information-frictions literature
(\citealt{stigler1961economics};
\citealt{genesove1997equity}; \citealt{anenberg2020housing}), is
that the Shen-Ross uniqueness premium is moderated by the renter
share of the housing stock; this primary hypothesis is supported in
the meta-regression with Hartung-Knapp coefficient
$\hat\beta_{\mathrm{renter}} = +0.025$ per standard deviation of
renter share, $95\%$ HKSJ CI $[+0.008, +0.042]$, two-sided
$p = 0.009$. This is direct evidence that the Shen-Ross uniqueness
effect is stronger in markets where a larger share of housing
transactions involves prospective tenants who lack the deep
neighborhood familiarity that long-term owner-occupiers possess.

The third reading combines the first two. Markets with high renter
shares have a true effect large enough that our per-city panel
detects it (San Francisco at $14.0\%$ per $\sigma$ in dollar terms,
Dallas at $4.9\%$). Markets with lower renter shares have a true
effect smaller than the minimum we can detect at our sample size, and
our nulls there are uninformative. We favor this reading as the most
parsimonious of the three, and we develop it in
Section~\ref{sec:meta-regression-12}.

\paragraph{Dollarization.} Converting the per-$\sigma$ point estimates
to per-swap dollar amounts at the median sale price in each market,
the Shen-Ross channel implies that a one-standard-deviation increase
in listing uniqueness corresponds to a median price difference of
$\$156{,}926$ in San Francisco, $\$20{,}720$ in Dallas, and $\$15{,}007$
in Seattle. The median-price normalization compresses the
substantively different impacts these effects would have on a buyer:
the Seattle figure represents $1.9\%$ of median price, the San
Francisco figure $14.0\%$, and the Dallas figure $4.9\%$.
